\section{把结果变成可以审核的成果}
\subsection{图文和执行证据各自核对}
正式制图实际比较了三目标3D、二维投影与平行坐标，区域V的地图小倍数与时间--区域矩阵，以及在线服务ECDF与条件分布。图型依据问题和数据选择：真实第三维才用3D；区域相邻移动需要空间关系；未匹配请求必须在统计分母中可见。不是把默认Notebook图换一种颜色就结束。

报告核验还检查图注是否对应真实样本、训练seed、验证/测试身份、波动区间和图轴。独立审稿发现，某些图只使用一个训练seed的五模拟日，不能被文字写成五训练seed；动态融合和静态前沿的预算也应分别注明。这些表述修复不重跑没有受影响的模型，但要重新构建真正的LaTeX源和PDF。

两本完成版Notebook在干净Kernel执行，保留教师Q/LT/Q-KM原断言，默认300回合没有被样本运行替代。完整参数扫描与固定结果另由可重算脚本、当前有效索引和输入/代码/产物哈希追溯。Notebook执行、重画图和报告编译是不同检查，任何一项不能冒充另外两项。

收尾时Codex把一个尚未解析命令行参数的评估入口当作支持\texttt{--help}的脚本调用，实际重放了冻结测试评估并覆盖部分文件。随后单独保存重放身份与可恢复检查点，独立Agent逐条确认95条episode/trace与原运行哈希相同，科学汇总及融合选择一致，仅运行耗时改变。未能恢复的旧逐日耗时字节如实登记；入口增加帮助参数和既存输出拒写检查。这个问题说明，命令行的无副作用行为也必须实际验证，不能靠脚本命名推断。

\subsection{公开发布边界也是一次真实反馈}
数据审计没有找到额外公开再分发许可。我把Codex的权限弹窗发给GPT询问，GPT建议采用“无额外授权，来源下载复现”，同时保留完整本地实验并限制公开数据类型。后续Codex交付记录表明，公开仓库采用了该数据边界。现有视频不能单独证明我曾逐条向Codex发出那段完整建议，因此不把GPT建议改写成我当时的逐项原话。原始签到、含数据的ZIP及可还原的逐条派生保持本地；这没有降低教师任务和实验范围。

独立包审查进一步实测：匿名的匹配选单矩阵与多条总目标可以形成满秩方程，恢复单笔fare。因此“删掉用户ID”还不够，真实EA的membership记录也从公开投影中移除，保留聚合前沿、收敛、参数、耗时、哈希和独立审核。原始完整结果没有被删除或改写。

原作者的下载渠道经过真实获取和校验。早期转换出现极少坐标末位差异，Research随后查明直接Decimal十位有效数字ROUND\_HALF\_UP和文件末尾无CRLF的规则，重新独立下载并转换后，两城市SHA与教师CSV逐字节一致。旧近似转换和失败诊断保留原身份；新输入重建仍准确标注作者来源，不倒写提供者，也不把输入一致扩大为任意环境随机产物bit一致。对数据许可的处理是发布范围控制，不是只剩总结而隐去实现。

\section{本轮保留的认识}
这次连续执行把我原先要求的“先保证正确”落实到具体对象：服务覆盖、相同目标的离线参考、共同HV、真实可行动作和物理移动成本。算法跑通以后，仍需检查比较是否成立、效果代价是什么；方法没有新增价值时，回退也构成有依据的结果。

我随后要求网页端GPT对固定提交进行二审，并在下一组真实截图中区分代码与科学成果、Process证据和最后提交的不同状态。
